
Since the explosion in the adoption of large language models (LLMs), researchers have sought to measure the growth in capabilities across a wide variety of benchmarks. This work has moved from measuring capabilities on abstract knowledge or capabilities \citep{hendrycks2021, srivastava2023} to more real-world, economically valuable tasks \citep{gdpval2025, profbench2025, occubench2026}. These benchmarks provide an excellent indication of a model's capacity to automate work. However, in many real organizational and personal settings, models do not produce the output themselves. They are often assisting other agents, such as a human or another, perhaps weaker, LLM to complete the task. In such settings, models help workers plan an approach, decompose a task into steps, identify potential errors, and revise their work while leaving responsibility for the final deliverable with the worker. A model that excels at completing a task independently may not be the model that provides the most useful guidance to someone else. Put differently, as research with humans shows, the best (AI) player may not be the best coach \citep{hinds1999, camerer1989}. 


% Yet, almost without exception, most benchmarks measure only a single dimension of the model as an autonomous task solver or how well a model completes a task on its own. In organizational settings, however, models are also used as assistants. They help workers plan an approach, decompose a task into steps, identify potential errors, and revise their work while leaving responsibility for the final deliverable with the worker. A model that excels at completing a task independently may not be the model that provides the most useful guidance to someone else. Put differently, the best player may not be the best coach.

%% we should also cite andy and erik's centaur benchpaper upfront. 
Existing research does not provide a general way to compare these two capabilities. On the one hand, large-scale computer science benchmarks primarily rank models as autonomous task solvers, including increasingly realistic evaluations of professional and economically valuable work \citep{hendrycks2021, liang2023, gdpval2025, sopbench2025, occubench2026, profbench2025, jobbench2026}. On the other hand, field experiments by economists and management scholars examine whether AI augments human performance, but generally evaluate a single model in a particular domain and against a limited set of experimental conditions \citep{noy2023, peng2023, dellacqua2023, chen2024, doshi2024, brynjolfsson2025, ju2026}. 

There is a line of work on multi-agent systems. In one stream, past work finds that a model's performance in isolation does not reliably predict its performance in interactive settings \citep{collins2024, chang2025, vaccaro2024}. Research has also studied how stronger models can teach weaker agents \citep{saha2023} or human learners \citep{macina2025}, and how AI can complement human decision makers \citep{bansal2021}. These studies evaluate assistance within a single domain (e.g. reasoning accuracy, tutoring pedagogy etc.), and score the assistance interaction itself. No existing benchmark compares models’ ability to assist and automate on an identical set of work-related tasks. Indeed, \citet{haupt2025} argue that model evaluation should move beyond an ``imitation game'' of autonomous task completion toward \emph{centaur evaluations} that measure how models augment human performance, yet no research operationalizes this call, jointly evaluating models as automators and as augmenters across multiple professional domains. Motivated by this gap, in this paper, we explore how models' automation capability correlates with its augmentation capability and explore how this correlation itself varies across tasks and models.

% Research has also studied how stronger models can teach weaker agents \citep{saha2023} or human learners \citep{macina2025}, and how AI can complement human decision makers \citep{bansal2021}. But, each paper examines a single assistant in one (narrow) domain and does not compare models against one another. 

%%This leaves organizations and individuals in a difficult position. When selecting an AI tool for human-in-the-loop work, ``Which model is best?'' is the wrong question; ``Best for what role, on what kind of task?'' is the right one.

We evaluate LLMs as automators and augmenters across seven heterogeneous professional tasks. Our pipeline evaluates a set of frontier LLMs across two axes simultaneously: usage mode (automation versus augmentation) and task domain. In automation mode, each focal model completes a professional task end-to-end. In augmentation mode, each focal model provides assistance text to a fixed worker model, which then completes the task. Holding the worker constant isolates the value of the assistant's guidance from the executor's own capability, producing a direct measure of assistant models quality rather than team outcome. This approach allows us to measure differences in augmentation and automation to any model-task pairing without a separate field experiment. Both outputs are evaluated by a panel of LLM judges using task-specific rubrics, with model families excluded from judging their own outputs (e.g. Claude Opus 4.8 cannot judge other Claude model outputs). The result is a rank matrix at the task and model level for both the augmentation and automation regimes. 

%Our broader framework has direct organizational implications. Concretely, the strongest direct solver is not the strongest assistant, and the best augmenting model changes from task to task, so a single leaderboard cannot tell an organization which model to deploy in which role. Rather than selecting models solely from general-purpose leaderboards, firms can use our pipeline to run task-specific evaluations that reflect their own workflows and deployment choices, whether it be delegating work entirely to a model, or keeping a human in the loop with model-generated scaffolding. More broadly, the results suggest that future benchmarks should move beyond agent-alone performance and incorporate role-specific, workflow-aware evaluations that test whether models create value as assistants, coaches, and collaborators, not only as independent task solvers.

Three key results are worth highlighting. First, we find that model rankings diverge meaningfully across usage modes and task domains. The strongest autonomous task solver is not necessarily the strongest assistant, and no single model dominates across augmentation tasks. The model-level rank correlation is about $0.48$ between augmentation and automation regimes and itself varies, ranging from $-0.04$ (travel planning) to $0.85$ (tax preparation) across tasks. Second, assistance is not always beneficial. While GPT-5-mini ranks first for both automation and augmentation schemes, the unaided worker baseline, in which GPT-3.5-Turbo completes the task without guidance from another model, achieves the second best average rank in the augmentation comparison. Guidance that is overly complex, poorly matched to the worker, or miscalibrated to the task can leave the downstream worker worse off than working alone. Third, reversals in model performance across regimes at the task level are large and stable. For example, the strongest automation model on our market trends analysis task, Claude-Opus-4.8 (mean rank 2.05), is among the weakest assistants on the same task (8.15), while GPT-4.1 shows the mirror-image profile on counseling (a lower-ranked direct solver but the strongest assistant).

%% 
Our study makes three contributions. First, we introduce a common evaluation framework for comparing models as both autonomous task solvers and assistants across professional domains. Second, by holding the downstream worker constant, we provide a direct measure of the marginal value of model-generated guidance and show that automation performance does not determine augmentation performance. Third, we demonstrate that the value of assistance depends on both the model and the task, challenging the use of a single general-purpose leaderboard for organizational model selection. These findings have direct implications for organizations choosing and deploying LLMs. The relevant question is not simply, ``Which model is best?'' but rather, ``Which model is best for this role and this task?'' A model selected to produce finished deliverables may differ from the model selected to support a worker through planning, critique, or revision. More broadly, our results suggest that future benchmarks should move beyond agent-alone performance and incorporate role-specific, workflow-aware evaluations of models as assistants, coaches, and collaborators.


%%% OLD
% Intro of the idea
% When an organization adopts a large language model (LLM), it faces two decisions: which model to use, and how to use it. For the first, the natural move is to consult a benchmark leaderboard, but almost without exception, leaderboards measure a single dimension, automation: how well a model completes a task on its own. This is rarely how models are used at work. Far more often a model augments work, assisting a person by planning the approach, breaking the task into steps, or prompting self-review while the worker produces the final deliverable. The gap between the automation that benchmarks reward and the augmentation that organizations actually deploy is what we examine in this paper. As we will show, the best player is not always the best coach.

% Related Works
% On one hand, computer scientists have built large-scale benchmarks to evaluate and rank how well LLMs complete tasks: MMLU \citep{hendrycks2021} tests knowledge across 57 subjects, HELM \citep{liang2023} evaluates 30 models on 42 scenarios, and a growing ecosystem of professional benchmarks now tests agents on everything from industrial workflows \citep{sopbench2025} to occupational tasks drawn from 65 specialized domains \citep{occubench2026}. Additionally, GDPval \citep{gdpval2025} and related benchmarks \citep{profbench2025, jobbench2026} have moved toward economically grounded tasks and pairwise evaluation. Yet \citet{collins2024} and \citet{chang2025} have shown that AI-alone accuracy does not predict interactive performance, and \citet{vaccaro2024}, in a meta-analysis of 106 experiments, find that human-AI combinations frequently underperform the best of either human alone or AI alone.

% On the other hand, economists and management scholars have run field experiments to measure AI's productivity effects: ChatGPT cut professional writing time by 40\% and raised quality by 18\% \citep{noy2023}, GitHub Copilot cut programming time by more than half \citep{peng2023}, and a large-scale deployment in customer support raised productivity by 15\% \citep{brynjolfsson2025}. Another issue is that field experiments that test augmentation are narrow by design: \citet{dellacqua2023} test one model with BCG consultants, \citet{chen2024} test ghostwriting versus sounding-board modes in advertising, \citet{doshi2024} study AI-assisted creative writing, \citet{ju2026} study human-AI versus human-human advertising teams. Each is informative but none tell you which of today's models is best suited to augmenting workers across a range of professional tasks---because none were designed to. They test a single model, in a single domain, against a single baseline. 

%This distinction matters beyond benchmarking: \citet{agrawal2023} argue that one worker's automation creates augmentation opportunities for another, and \citet{marguerit2025} finds that augmentation AI raises wages for high-skilled workers while automation AI displaces low-skilled ones. If benchmarks only measure what a model can do alone, and field experiments only measure whether a specific tool helped in a specific setting, we are missing a big part of the picture. Both fields have produced genuine insight but neither fully addresses how workers and organizations actually want to use AI. In practice, a model rarely produces a finished deliverable on its own. Instead it assists another individual. For example, a nationwide survey of 1,500 U.S. workers across 104 occupations found that ``equal partnership'' with AI, not replacement, was the dominant preference in nearly half of all occupational categories \citep{shao2025}.

% Introducing the gap
% All in all, no existing framework jointly evaluates models as automators and as augmenters across multiple professional domain; the cross-model, cross-mode comparison that would actually inform tool selection does not exist. Therefore, there is a gap that captures the distinction between \emph{automation} and \emph{augmentation}. A model that performs well when asked to complete a task end-to-end may not be the same model that best helps a human worker complete that same task. The best players are not always the best coaches.  This leaves organizations and individuals in a difficult position. When selecting an AI tool for human-in-the-loop work, ``Which model is best?'' is the wrong question; ``Best for what role, on what kind of task?'' is the right one. \citet{haupt2025} make the methodological implication explicit: evaluation-by-imitation rewards systems that replace people rather than systems that augment them, and the field needs evaluations where humans and AI solve tasks together. 

% Two close papers to ours
% \citet{fugener2026} further formalize this distinction and show that the optimal mode depends on the type of complementarity between human and AI capability. Industry has begun to move in this direction as well. The closest prior effort is Upwork's Human+Agent Productivity Index \citep{upwork_hapi2025}, which evaluates frontier agents on roughly 300 real, paid client projects and finds that pairing agents with expert freelancers raises completion rates by up to 70\% relative to agents alone. \citep{upwork_hapi2025} is important evidence that agent-only scoring understates economic value, but it differs from our study in three ways. First, it studies a single human--AI pairing rather than augmentation \emph{across} models: it contrasts agent-alone with human-in-the-loop, not model A versus model B as assistants to a fixed worker. Second, in HAPI a human expert assists the AI, reviewing and correcting the agent's output, so the human occupies the coaching role and the model is the party being helped. Our framework complements theirs. Instead of treating the models as direct task solvers, we regard them as the assistant, providing scaffolding to a fixed worker, so we evaluate models in the coaching role rather than as the worker being coached. This is precisely the capability our analysis foregrounds---whether a model is a good coach, not only a good player. Third, HAPI's coverage is bounded by what can be sourced and staffed on a freelance marketplace and each task requires recruiting a qualified expert, while our fixed-worker design scales to any professional task we can specify, without freelancer recruitment. HAPI thus motivates a broader evaluation gap: existing benchmarks largely measure agent-alone performance, while we study which models generate the greatest marginal value as assistants to a fixed worker.

%\textcolor{red}{Kris: I still feel like we're doing too much lit review before we get to the real point}

% What we add
% To our knowledge, we are the first to evaluate LLMs jointly as automators and augmenters across multiple professional domains. Our pipeline evaluates a set of frontier LLMs across two axes simultaneously: usage mode (automation versus augmentation) and task domain. In automation mode, each model completes a professional task end-to-end. In augmentation mode, each model provides coaching to a fixed worker model, which then completes the task. Holding the worker constant isolates the value of the assistant's guidance from the executor's own capability, producing a direct measure of coaching quality rather than team outcome. This method uses GPT 3.5 Turbo to simulate a worker, which allows us to approximate centaur evaluations. This approach allows us to measure differences in augmentation and automation to any model-task pairing without a separate field experiment.

%Both outputs are evaluated by a panel of LLM judges using task-specific rubrics across seven heterogeneous professional domains, with model families excluded from judging their own outputs (e.g. Claude Opus 4.8 cannot judge other Claude model outputs). The result is not only a single leaderboard score but also a cross-regime rank matrix: for each model, we can ask whether it ranks higher as a direct solver or as a planning coach, and whether that ordering holds across tasks. Interestingly, the unaided worker baseline, in which GPT-3.5-Turbo completes the task without external assistance, achieves the best average augmentation rank overall. This suggests that assistance is not automatically beneficial and poorly matched, overly complex, or miscalibrated guidance can leave a downstream worker worse off than working alone. Overall, we find model rankings diverge meaningfully across modes and across task domains---the best automator is not always the best augmenter, and no single model dominates across augmentation tasks. 

%Our broader framework has direct organizational implications. Concretely, the strongest direct solver is not the strongest assistant, and the best augmenting model changes from task to task, so a single leaderboard cannot tell an organization which model to deploy in which role. Rather than selecting models solely from general-purpose leaderboards, firms can use our pipeline to run task-specific evaluations that reflect their own workflows and deployment choices, whether it be delegating work entirely to a model, or keeping a human in the loop with model-generated scaffolding. More broadly, the results suggest that future benchmarks should move beyond agent-alone performance and incorporate role-specific, workflow-aware evaluations that test whether models create value as assistants, coaches, and collaborators, not only as independent task solvers.
